problem:  
Let \((M^{n+1}, g)\) be a globally hyperbolic Lorentzian manifold obeying the **null energy condition**, and let \(\Sigma \subset M\) be a smooth, spacelike Cauchy hypersurface. Suppose that within \(\Sigma\) there exists a compact, two‑sided **stable marginally outer trapped surface** \(S \subset \Sigma\), whose stability operator \(\mathcal L_S\) has lowest eigenvalue \(\lambda_1(\mathcal L_S) = 0\).  

Let \(\ell_+\) be the future‑directed null normal to \(S\) satisfying \(\theta_+(S)=0\). Define the associated **null second fundamental form** \(\chi_+\), with tracefree part \(\sigma_+\), and let \(\omega\) be the connection one‑form of the normal bundle. Assume \(S\) has positive Yamabe type, i.e. the conformal class of its induced metric admits metrics of positive scalar curvature.

**Problem:**  
Does the stability marginality \(\lambda_1(\mathcal L_S)=0\) together with positive Yamabe type of \(S\) force the Ricci tensor to satisfy a *null splitting* condition along the null hypersurface \(\mathcal N^+(S)\) generated by \(\ell_+\)?  In particular, is it true (perhaps under additional bounded curvature assumptions) that \(\mathcal N^+(S)\) must be a **totally geodesic, non‑expanding null hypersurface** with \(\mathrm{Ric}(\ell_+, \ell_+) = 0\), implying that a neighborhood of \(S\) locally splits as a warped product
\[
([0,\varepsilon) \times S,\; -2\,du\,dt + h_S),
\]
where \(h_S\) is the induced metric on \(S\)?  

Equivalently: is there a general **null splitting rigidity theorem** stating that, if a stable marginally outer trapped surface of positive Yamabe type saturates stability (\(\lambda_1=0\)) under the null energy condition, then the ambient spacetime is locally “optically flat” along the associated congruence?

---

Why is it a "good" problem:  
1. **Profound insight:**  
This asks whether stability saturation of a MOTS—a quasi‑local condition—forces *local null geometric rigidity* of the ambient spacetime. It distills a conjectural equivalence between an *eigenvalue extremality* (analytical condition) and *null splitting* (geometric condition), probing the borderline between formation of trapped surfaces and rigidity of isolated horizon structures.

2. **Bridge between fields:**  
It intertwines analysis (spectral theory of the MOTS stability operator) with Lorentzian comparison geometry and causal analysis. A solution would fuse PDE‑based methods from geometric analysis with techniques of null congruences and Raychaudhuri evolution known in general relativity.

3. **Extreme simplicity:**  
The question reads almost elementarily—*does equality in a stability condition force null rigidity?*—but encompasses deep interplay among curvature inequalities, stability spectra, and causal geometry. The ingredients are classical, the connection among them completely unexplored.

4. **Potential to open new directions:**  
An affirmative answer would establish a Lorentzian **null splitting theorem** paralleling the Cheeger–Gromoll style rigidity in a quasi‑local, null context. A counterexample would clarify the boundaries of marginal stability and help classify near‑horizon geometries analytically. Either outcome could generate a new chapter in Lorentzian comparison geometry focused on null optical structures.

By combining analytic extremality (λ₁=0), quasi‑local curvature type, and global energy conditions in a minimal and precise way, this problem attains the balance of clarity, challenge, and conceptual depth that defines a “good’’ research problem.